\textbf{1. Kiến thức cần nhớ}
\begin{itemize}
\item \textbf{Phân tích đa thức thành nhân tử} (hay \textit{thừa số}) là biến đổi đa thức đó thành một tích của những đa thức.
\item Các phương pháp thường dùng:
\begin{itemize}
\item \textbf{Đặt nhân tử chung}: $AB+AC=A(B+C)$.
\item \textbf{Dùng hằng đẳng thức}: bảy hằng đẳng thức đáng nhớ (bình phương, hiệu hai bình phương, lập phương, tổng và hiệu hai lập phương).
\item \textbf{Nhóm hạng tử}: nhóm thích hợp để mỗi nhóm có nhân tử chung hoặc là một hằng đẳng thức.
\end{itemize}
\item Đổi dấu khi cần: $A-B=-(B-A)$.
\end{itemize}

\textbf{2. Phương pháp giải các dạng}

%% DANG 8D2?9-1
\textbf{Dạng 1. Đặt nhân tử chung}
\begin{enumerate}
\item Tìm nhân tử chung: ƯCLN của các hệ số và các biến (lấy số mũ nhỏ nhất) hoặc đa thức chung.
\item Viết đa thức thành tích của nhân tử chung với đa thức còn lại; kiểm tra bằng cách nhân lại.
\end{enumerate}
Lưu ý: đổi dấu để làm xuất hiện nhân tử chung, ví dụ $x(y-1)+3(1-y)=(y-1)(x-3)$.

%% DANG 8D2?9-2
\textbf{Dạng 2. Dùng hằng đẳng thức}
\begin{enumerate}
\item Đếm số hạng tử: 2 hạng tử $\to$ thử $A^2-B^2$, $A^3\pm B^3$; 3 hạng tử $\to$ thử $(A\pm B)^2$; 4 hạng tử $\to$ thử $(A\pm B)^3$.
\item Viết các hạng tử dưới dạng $A^2$, $B^2$, $A^3$, $B^3$ rồi áp dụng hằng đẳng thức.
\end{enumerate}

%% DANG 8D2?9-3
\textbf{Dạng 3. Nhóm hạng tử, phối hợp nhiều phương pháp}
\begin{enumerate}
\item Đặt nhân tử chung trước (nếu có).
\item Nhóm các hạng tử để mỗi nhóm có nhân tử chung hoặc là hằng đẳng thức (ví dụ $x^2-2x+1-y^2=(x-1)^2-y^2$).
\item Tiếp tục phân tích đến khi không phân tích được nữa.
\end{enumerate}
Lưu ý: nhóm đúng thì sau khi phân tích từng nhóm, các nhóm phải còn nhân tử chung.

%% DANG 8D2?9-4
\textbf{Dạng 4. Ứng dụng: tìm x, tính nhanh, chứng minh chia hết}
\begin{enumerate}
\item Tìm $x$: chuyển hết về một vế, phân tích thành tích, dùng $A\cdot B=0\iff A=0$ hoặc $B=0$.
\item Tính nhanh: phân tích thành tích có thừa số tròn, ví dụ $37\cdot 63+37\cdot 37=37\cdot 100$.
\item Chứng minh chia hết: phân tích thành tích có một thừa số chia hết cho số cần chứng minh.
\end{enumerate}
Lưu ý: khi tìm $x$ không chia hai vế cho biểu thức chứa $x$ (dễ mất nghiệm).
